\documentclass[11pt, a4paper]{problems_template}

\begin{document}

\begin{titlepage}
  \begin{center}
	\Huge{Прасолов - Задачи по Планиметрии}
  \end{center}
\end{titlepage}   

\chapter{Подобные треугольники}

\section{Вводные задачи}

\begin{exercise}
В остроугольном треугольнике $ABC$ проведены высоты $AA_{1}$ и $BB_{1}$. Докажите, что $A_{1} C \cdot BC = B_{1}C \cdot AC$.
\end{exercise}

\begin{Solution}
Треугольники $AA_{1}C$ и $BB_{1}C$ подобны по прямому углу и углу при вершине $C$. Значит отношение гипотенуз этих треугольников равно отношению соответствующих катетов при вершине $C$. 
$$
\frac{AC}{BC} = \frac{A_{1}C}{B_{1}C}
$$
Откуда следует искомое равенство.
\end{Solution}

\begin{exercise}
Докажите, что медианы треугольника пересекаются в одной точке и делятся этой точкой в отношении $2 : 1$, считая от вершины.
\end{exercise}

\begin{Solution}
Рассмотрим треугольник $ABC$ и две медианы $AA_{1}$ и $BB_{1}$, которые пересекаются в точке $O$. Тогда $A_{1}B_{1}$ будет средней линией $ABC$. Значит
$$
A_{1}B_{1} \parallel AB \implies \angle A_{1}B_{1}O = \angle ABO, \quad \angle B_{1}A_{1}O = \angle BAO 
$$
Значит треугольник $A_{1}B_{1}O$ подобен треугольнику $ABO$ с коэффициентом подобия $k = 2$. И точка $O$ делит обе медианы в отношении $2 : 1$, считая от вершины. Аналогичные соображения можно применить к медианам $BB_{1}$ и $CC_{1}$. Отсюда следует утверждение задачи.
\end{Solution}

\section{Отрезки, заключённые между параллельными прямыми}

\begin{exercise}
Через точку $P$ медианы $CC_{1}$ треугольника $ABC$ проведены прямые $AA_{1}$ и $BB_{1}$ (точки $A_{1}$ и $B_{1}$ лежат на сторонах $BC$ и $CA$). Докажите, что $A_{1} B_{1} \parallel AB$.
\end{exercise}



\end{document} 
